\subsection{Операторнозначные функции комплексного переменного. Аналитичность резольвенты. Спектральный радиус. Основная теорема о спектре}

\begin{definition}[Операторнозначная функция комплексного переменного]
	Пусть \(E\) --- банахово пространство над \(\mathbb{C}\), а \(D\subset\mathbb{C}\).
	Отображение
	\[
	F:D\to\mathcal{L}(E)
	\]
	называется \textbf{операторнозначной функцией комплексного переменного}.
\end{definition}

\begin{definition}[Аналитичность операторнозначной функции]
	Функция
	\[
	F:D\to\mathcal{L}(E)
	\]
	называется \textbf{аналитической} в точке \(\lambda_0\in D\), если существует предел
	\[
	F'(\lambda_0)
	=
	\lim_{\lambda\to\lambda_0}
	\frac{F(\lambda)-F(\lambda_0)}{\lambda-\lambda_0}
	\]
	в норме пространства \(\mathcal{L}(E)\).
	
	Если \(F\) аналитична в каждой точке области \(D\), то \(F\) называется аналитической в \(D\).
\end{definition}


\begin{definition}[Резольвентное множество, резольвента и спектр]
	Пусть \(E\) --- комплексное банахово пространство,
	\[
	A\in\mathcal{L}(E).
	\]
	Для \(\lambda\in\mathbb{C}\) обозначим
	\[
	A_\lambda=A-\lambda I.
	\]
	
	Число \(\lambda\in\mathbb{C}\) называется \textbf{регулярным значением} оператора \(A\), если оператор
	\[
	A_\lambda=A-\lambda I
	\]
	непрерывно обратим, то есть
	\[
	A_\lambda^{-1}\in\mathcal{L}(E).
	\]
	
	\textbf{Резольвентным множеством} называется множество
	\[
	\rho(A)=\{\lambda\in\mathbb{C}: A_\lambda^{-1}\in\mathcal{L}(E)\}.
	\]
	
	Для \(\lambda\in\rho(A)\) оператор
	\[
	R_\lambda=R_\lambda(A)=(A-\lambda I)^{-1}
	\]
	называется \textbf{резольвентой} оператора \(A\).
	
	\textbf{Спектром} оператора \(A\) называется множество
	\[
	\sigma(A)=\mathbb{C}\setminus\rho(A).
	\]
\end{definition}

\begin{definition}[Спектральный радиус]
	Спектральным радиусом оператора \(A\) называется число
	\[
	r(A)=\sup_{\lambda\in\sigma(A)}|\lambda|.
	\]
\end{definition}

\begin{proposition}[Резольвента вне круга радиуса \(\|A\|\)]
	Если
	\[
	|\lambda|>\|A\|,
	\]
	то
	\[
	\lambda\in\rho(A),
	\]
	причём
	\[
	R_\lambda
	=
	-\frac{1}{\lambda}
	\sum_{n=0}^{\infty}\frac{A^n}{\lambda^n}.
	\]
\end{proposition}

\begin{proof}
	\textbf{1. Сведение к ряду Неймана.}
	
	Имеем
	\[
	A-\lambda I
	=
	-\lambda\left(I-\frac{1}{\lambda}A\right).
	\]
	Так как
	\[
	\left\|\frac{1}{\lambda}A\right\|
	=
	\frac{\|A\|}{|\lambda|}
	<1,
	\]
	то по теореме о ряде Неймана оператор
	\[
	I-\frac{1}{\lambda}A
	\]
	обратим, и
	\[
	\left(I-\frac{1}{\lambda}A\right)^{-1}
	=
	\sum_{n=0}^{\infty}\frac{A^n}{\lambda^n}.
	\]
	
	\textbf{2. Формула для резольвенты.}
	
	Из разложения
	\[
	A-\lambda I
	=
	-\lambda\left(I-\frac{1}{\lambda}A\right)
	\]
	получаем
	\[
	(A-\lambda I)^{-1}
	=
	-\frac{1}{\lambda}
	\left(I-\frac{1}{\lambda}A\right)^{-1}.
	\]
	Следовательно,
	\[
	R_\lambda
	=
	-\frac{1}{\lambda}
	\sum_{n=0}^{\infty}\frac{A^n}{\lambda^n}.
	\]
	Значит,
	\[
	\lambda\in\rho(A).
	\]
\end{proof}

\begin{corollary}
	Спектр оператора \(A\) лежит в замкнутом круге:
	\[
	\sigma(A)\subset\{\lambda\in\mathbb{C}:|\lambda|\leqslant \|A\|\}.
	\]
\end{corollary}

\begin{theorem}[Аналитичность резольвенты]
	Резольвентное множество \(\rho(A)\) открыто, а функция
	\[
	R_\lambda=(A-\lambda I)^{-1}
	\]
	аналитична на \(\rho(A)\). Более того,
	\[
	R'_\lambda=R_\lambda^2.
	\]
\end{theorem}

\begin{proof}
	\textbf{1. Открытость резольвентного множества.}
	
	Возьмём произвольную точку
	\[
	\lambda_0\in\rho(A).
	\]
	Тогда оператор
	\[
	R_{\lambda_0}=(A-\lambda_0I)^{-1}
	\]
	существует и ограничен.
	
	Пусть \(\lambda\) близко к \(\lambda_0\). Обозначим
	\[
	\delta=\lambda-\lambda_0.
	\]
	Тогда
	\[
	A-\lambda I
	=
	A-\lambda_0I-\delta I
	=
	(A-\lambda_0I)(I-\delta R_{\lambda_0}).
	\]
	Если
	\[
	|\delta|\mathbb{C}dot\|R_{\lambda_0}\|<1,
	\]
	то по теореме о ряде Неймана оператор
	\[
	I-\delta R_{\lambda_0}
	\]
	обратим. Значит, обратим и оператор \(A-\lambda I\).
	
	Следовательно, при
	\[
	|\lambda-\lambda_0|<\frac{1}{\|R_{\lambda_0}\|}
	\]
	имеем
	\[
	\lambda\in\rho(A).
	\]
	Значит, \(\rho(A)\) открыто.
	
	\textbf{2. Локальное разложение резольвенты.}
	
	Из равенства
	\[
	A-\lambda I
	=
	(A-\lambda_0I)(I-(\lambda-\lambda_0)R_{\lambda_0})
	\]
	получаем
	\[
	R_\lambda
	=
	(I-(\lambda-\lambda_0)R_{\lambda_0})^{-1}R_{\lambda_0}.
	\]
	При
	\[
	|\lambda-\lambda_0|\mathbb{C}dot\|R_{\lambda_0}\|<1
	\]
	по ряду Неймана:
	\[
	(I-(\lambda-\lambda_0)R_{\lambda_0})^{-1}
	=
	\sum_{n=0}^{\infty}(\lambda-\lambda_0)^nR_{\lambda_0}^n.
	\]
	Следовательно,
	\[
	R_\lambda
	=
	\sum_{n=0}^{\infty}(\lambda-\lambda_0)^nR_{\lambda_0}^{n+1}.
	\]
	Это степенной ряд по \(\lambda-\lambda_0\) со значениями в \(\mathcal{L}(E)\). Поэтому
	\[
	\lambda\mapsto R_\lambda
	\]
	аналитична в окрестности \(\lambda_0\).
	
	\textbf{3. Формула для производной.}
	
	Из локального разложения сразу получаем
	\[
	R'_{\lambda_0}=R_{\lambda_0}^2.
	\]
	Так как \(\lambda_0\in\rho(A)\) была произвольной,
	\[
	R'_\lambda=R_\lambda^2
	\qquad
	\forall\lambda\in\rho(A).
	\]
\end{proof}

\begin{remark}[Равенство Гильберта для резольвенты]
	Если
	\[
	\lambda,\mu\in\rho(A),
	\]
	то
	\[
	R_\lambda-R_\mu
	=
	(\lambda-\mu)R_\lambda R_\mu.
	\]
	
	Действительно,
	\[
	R_\lambda-R_\mu
	=
	R_\lambda R_\mu^{-1}R_\mu
	-
	R_\lambda R_\lambda^{-1}R_\mu
	=
	R_\lambda(R_\mu^{-1}-R_\lambda^{-1})R_\mu.
	\]
	Но
	\[
	R_\mu^{-1}=A-\mu I,
	\qquad
	R_\lambda^{-1}=A-\lambda I.
	\]
	Значит,
	\[
	R_\mu^{-1}-R_\lambda^{-1}
	=
	(A-\mu I)-(A-\lambda I)
	=
	(\lambda-\mu)I.
	\]
	Следовательно,
	\[
	R_\lambda-R_\mu
	=
	(\lambda-\mu)R_\lambda R_\mu.
	\]
\end{remark}

\begin{theorem}[Основная теорема о спектре]
	Пусть \(E\) --- комплексное банахово пространство и
	\[
	A\in\mathcal{L}(E).
	\]
	Тогда:
	\begin{itemize}
		\item спектр \(\sigma(A)\) --- непустой компакт;
		\item спектральный радиус задаётся формулой
		\[
		r(A)
		=
		\lim_{n\to\infty}\sqrt[n]{\|A^n\|}.
		\]
	\end{itemize}
\end{theorem}

\begin{proof}
	\textbf{1. Замкнутость и ограниченность спектра.}
	
	Из предыдущего утверждения знаем, что
	\[
	\rho(A)
	\]
	открыто. Поэтому его дополнение
	\[
	\sigma(A)=\mathbb{C}\setminus\rho(A)
	\]
	замкнуто.
	
	Кроме того, если
	\[
	|\lambda|>\|A\|,
	\]
	то
	\[
	\lambda\in\rho(A).
	\]
	Значит,
	\[
	\sigma(A)\subset\{\lambda\in\mathbb{C}:|\lambda|\leqslant\|A\|\}.
	\]
	Следовательно, \(\sigma(A)\) замкнуто и ограничено. Если оно непусто, то оно компактно.
	
	Осталось доказать непустоту.
	
	\textbf{2. Непустота спектра.}
	
	Докажем от противного. Пусть
	\[
	\sigma(A)=\varnothing.
	\]
	Тогда
	\[
	\rho(A)=\mathbb{C},
	\]
	то есть резольвента
	\[
	R_\lambda=(A-\lambda I)^{-1}
	\]
	определена и аналитична на всей комплексной плоскости.
	
	При \(|\lambda|>\|A\|\) уже есть формула
	\[
	R_\lambda
	=
	-\frac{1}{\lambda}
	\sum_{n=0}^{\infty}\frac{A^n}{\lambda^n}.
	\]
	Поэтому
	\[
	\|R_\lambda\|
	\leqslant
	\frac{1}{|\lambda|}
	\sum_{n=0}^{\infty}
	\frac{\|A\|^n}{|\lambda|^n}
	=
	\frac{1}{|\lambda|}
	\mathbb{C}dot
	\frac{1}{1-\frac{\|A\|}{|\lambda|}}
	=
	\frac{1}{|\lambda|-\|A\|}.
	\]
	Отсюда
	\[
	\|R_\lambda\|\to0
	\qquad
	\text{при }|\lambda|\to\infty.
	\]
	
	Значит, \(R_\lambda\) ограничена вне большого круга. Внутри этого круга она ограничена по
	непрерывности на компакте. Следовательно, \(R_\lambda\) ограничена на всей \(\mathbb{C}\).
	
	По теореме Лиувилля для операторнозначных функций \(R_\lambda\) постоянна. Так как
	\[
	R_\lambda\to0
	\qquad
	|\lambda|\to\infty,
	\]
	эта постоянная равна нулю:
	\[
	R_\lambda\equiv0.
	\]
	Но \(R_\lambda\) является обратным оператором к \(A-\lambda I\), значит не может быть нулевым
	оператором. Получили противоречие.
	
	Следовательно,
	\[
	\sigma(A)\ne\varnothing.
	\]
	Значит, \(\sigma(A)\) --- непустой компакт.
	
	\textbf{3. Формула спектрального радиуса.}
	
	Обозначим
	\[
	L=\lim_{n\to\infty}\sqrt[n]{\|A^n\|}.
	\]
	Предел существует по стандартной лемме для субмультипликативной последовательности,
	поскольку
	\[
	\|A^{n+m}\|\leqslant\|A^n\|\mathbb{C}dot\|A^m\|.
	\]
	
	Докажем, что
	\[
	r(A)=L.
	\]
	
	\begin{itemize}
		\item \textbf{Оценка \(r(A)\leqslant L\).}
		
		Пусть
		\[
		\lambda\in\sigma(A).
		\]
		Тогда для любого \(n\in\mathbb{N}\)
		\[
		\lambda^n\in\sigma(A^n).
		\]
		Действительно, если бы \(A^n-\lambda^nI\) был обратим, то из разложения
		\[
		A^n-\lambda^nI
		=
		(A-\lambda I)
		(A^{n-1}+\lambda A^{n-2}+\ldots+\lambda^{n-1}I)
		\]
		и коммутативности множителей следовала бы обратимость \(A-\lambda I\), что противоречит
		\[
		\lambda\in\sigma(A).
		\]
		
		Так как спектр любого ограниченного оператора лежит в круге радиуса его нормы, то
		\[
		|\lambda|^n
		=
		|\lambda^n|
		\leqslant
		\|A^n\|.
		\]
		Следовательно,
		\[
		|\lambda|
		\leqslant
		\sqrt[n]{\|A^n\|}
		\qquad
		\forall n.
		\]
		Переходя к пределу по \(n\), получаем
		\[
		|\lambda|\leqslant L.
		\]
		Так как \(\lambda\in\sigma(A)\) произвольно,
		\[
		r(A)\leqslant L.
		\]
		
		\item \textbf{Оценка \(L\leqslant r(A)\).}
		
		Для \(|\lambda|>\|A\|\) мы знаем разложение
		\[
		R_\lambda
		=
		-\frac{1}{\lambda}
		\sum_{n=0}^{\infty}
		\frac{A^n}{\lambda^n}.
		\]
		Это ряд Лорана резольвенты в окрестности бесконечности.
		
		Так как резольвента аналитична на \(\rho(A)\), а
		\[
		\sigma(A)\subset\{\lambda:|\lambda|\leqslant r(A)\},
		\]
		то \(R_\lambda\) аналитична во всей внешности круга
		\[
		|\lambda|>r(A).
		\]
		По единственности ряда Лорана указанное разложение является лорановским разложением
		резольвенты во всей этой внешней области.
		
		Значит, ряд
		\[
		\sum_{n=0}^{\infty}
		\frac{A^n}{\lambda^{n+1}}
		\]
		сходится при каждом
		\[
		|\lambda|>r(A).
		\]
		Возьмём любое число \(s>r(A)\) и положим \(|\lambda|=s\). Из сходимости ряда следует, что
		его общий член стремится к нулю:
		\[
		\left\|
		\frac{A^n}{\lambda^{n+1}}
		\right\|
		=
		\frac{\|A^n\|}{s^{n+1}}
		\to0.
		\]
		Отсюда
		\[
		\limsup_{n\to\infty}\sqrt[n]{\frac{\|A^n\|}{s^{n+1}}}
		\leqslant1.
		\]
		То есть
		\[
		L\leqslant s.
		\]
		Так как это верно для любого \(s>r(A)\), получаем
		\[
		L\leqslant r(A).
		\]
	\end{itemize}
	
	Из двух оценок следует
	\[
	r(A)=L.
	\]
	То есть
	\[
	r(A)
	=
	\lim_{n\to\infty}\sqrt[n]{\|A^n\|}.
	\]
\end{proof}

\begin{idea}[Скелет билета]
	\begin{itemize}
		\item Для \(\lambda\in\mathbb{C}\) рассматриваем
		\[
		A_\lambda=A-\lambda I.
		\]
		
		\item Резольвентное множество:
		\[
		\rho(A)=\{\lambda: A_\lambda^{-1}\in\mathcal{L}(E)\}.
		\]
		
		\item Резольвента:
		\[
		R_\lambda=(A-\lambda I)^{-1}.
		\]
		
		\item Спектр:
		\[
		\sigma(A)=\mathbb{C}\setminus\rho(A).
		\]
		
		\item Если \(|\lambda|>\|A\|\), то
		\[
		R_\lambda
		=
		-\frac{1}{\lambda}
		\sum_{n=0}^{\infty}\frac{A^n}{\lambda^n}.
		\]
		Значит,
		\[
		\sigma(A)\subset\{|\lambda|\leqslant\|A\|\}.
		\]
		
		\item \(\rho(A)\) открыто, потому что обратимость сохраняется при малом возмущении:
		\[
		A-\lambda I
		=
		(A-\lambda_0I)(I-(\lambda-\lambda_0)R_{\lambda_0}).
		\]
		
		\item Резольвента аналитична, потому что локально раскладывается в ряд Неймана:
		\[
		R_\lambda
		=
		\sum_{n=0}^{\infty}
		(\lambda-\lambda_0)^nR_{\lambda_0}^{n+1}.
		\]
		
		\item Основная теорема:
		\[
		\sigma(A)\text{ --- непустой компакт},
		\qquad
		r(A)=\lim_{n\to\infty}\sqrt[n]{\|A^n\|}.
		\]
	\end{itemize}
\end{idea}